
\documentclass[12pt, a4paper]{article}
\usepackage[utf8]{inputenc}
\usepackage[italian]{babel}
\usepackage{graphicx}
\usepackage{listings}
\usepackage{xcolor}
\usepackage{hyperref}

% Configurazione per il codice Java
\definecolor{codegreen}{rgb}{0,0.6,0}
\definecolor{codegray}{rgb}{0.5,0.5,0.5}
\definecolor{codepurple}{rgb}{0.58,0,0.82}
\definecolor{backcolour}{rgb}{0.95,0.95,0.92}

\lstset{
    backgroundcolor=\color{backcolour},
    commentstyle=\color{codegreen},
    keywordstyle=\color{magenta},
    numberstyle=\tiny\color{codegray},
    stringstyle=\color{codepurple},
    basicstyle=\ttfamily\footnotesize,
    breakatwhitespace=false,
    breaklines=true,
    captionpos=b,
    keepspaces=true,
    numbers=left,
    numbersep=5pt,
    showspaces=false,
    showstringspaces=false,
    showtabs=false,
    tabsize=2,
    language=Java
}

\title{Relazione Progetto: Sistema Gestione Biblioteca}
\author{Studente: Marco}
\date{Maggio 2026}

\begin{document}

\maketitle

\tableofcontents
\newpage

\section{Introduzione}
Il presente progetto riguarda la realizzazione di un sistema per la gestione dei prestiti e della ricerca di libri in una biblioteca universitaria. Lo sviluppo segue i principi dell'Ingegneria del Software, focalizzandosi sulla tracciabilità dei requisiti, la modularità del codice e la verifica tramite test automatici.

\section{Analisi dei Requisiti}
\subsection{Casi d'Uso Principali}
Il sistema implementa due funzionalità core, fondamentali per il flusso di lavoro della biblioteca:
\begin{itemize}
    \item \textbf{UC1: Ricerca Libro:} Permette all'utente di verificare la presenza di un titolo nel catalogo mediante una ricerca testuale.
    \item \textbf{UC2: Gestione Prestito:} Gestisce la transazione di prestito verificando preventivamente la disponibilità del volume e aggiornandone lo stato.
\end{itemize}

\section{Architettura del Sistema}
\subsection{Modellazione dei Dati}
L'architettura segue il principio di Singola Responsabilità (SRP). I modelli (\texttt{Libro}, \texttt{Utente}) sono separati dalla logica di gestione (\texttt{Catalogo}, \texttt{GestorePrestiti}).

\subsection{Configurazione Maven}
Il progetto utilizza Maven come strumento di build e gestione delle dipendenze. Il file \texttt{pom.xml} è configurato per utilizzare \textbf{Java 21} e include le librerie \textbf{JUnit 5} per il testing.

\section{Verifica e Test}
\subsection{Test Unitari (JUnit)}
Per garantire la robustezza del sistema, è stato implementato un test unitario automatizzato per il flusso di prestito. Il test verifica che il sistema neghi correttamente il prestito di un libro già impegnato.

\begin{lstlisting}[caption={Esempio di Test Unitario per il Prestito}]
@Test
void testPrestitoLibroGiaOccupato() {
    gestore.registraPrestito(utente1, libro);
    boolean risultato = gestore.registraPrestito(utente2, libro);
    assertFalse(risultato);
}
\end{lstlisting}

\subsection{Risultati della Compilazione}
L'esecuzione del comando \texttt{mvn test} ha prodotto i seguenti risultati, confermando la correttezza della logica implementata:
\begin{itemize}
    \item \textbf{Build Status:} SUCCESS
    \item \textbf{Tests run:} 1
    \item \textbf{Failures:} 0
    \item \textbf{Errors:} 0
\end{itemize}

\section{Conclusioni}
Il sistema soddisfa i requisiti minimi richiesti, dimostrando una corretta integrazione tra modellazione UML, implementazione Java 21 e verifica tramite pipeline di test.

\end{document}
